% HD-F1-T22. GPT-6 Astra Pro; /root/transport; 22 November 2026 UTC.
\section[HD-F1-T22: ordinary alternating layers over the seven non-Gaussian fields]{HD-F1-T22: the original nonreal coefficient field excludes every ordinary alternating layer over the seven non-Gaussian fields}
\label{sec:f1-ordinary-alternating-source}

Retain the actual natural sixfold source, coefficient representation
and selected primitive rational corner of T13--T20. For
\[
 K=\mathbf Q(\sqrt{-d}),\qquad
 d\in\{2,7,11,19,43,67,163\},
 \tag{F1T22.1}
\]
the remaining coefficient image cannot have almost simple
projective image with socle \(A_n\), \(n\ge7\), and
actual perfect layer isomorphic to \(A_n\).
Here ``actual perfect layer'' means the commutator subgroup
\(L=P'\) of the inverse image \(P\) of this socle in
the original matrix group \(G\). This statement is about
\(L\simeq A_n\). It does not identify a double or triple
central cover with \(A_n\). The Gaussian field is also
outside the present assertion.

Write \(V=K^m\) for the original faithful absolutely
irreducible coefficient module. Its original inertia tuple
has nonidentity entries, generates \(G\), and satisfies
\[
 g_1\cdots g_b=1,\qquad b\ge12,\qquad m\ge27,
 \qquad\sum_{j=1}^b\operatorname{rank}(g_j-1)=2m+6.
 \tag{F1T22.2}
\]
These are the same period and Chevalley--Weil conditions
used in T13. We prove below that every nonidentity
\(g\in G\) would have moving rank at least \(m/5\).
Equation (F1T22.2) would then imply \(m\le15\).

\paragraph{1. The exact character field comes from the original corner.}
The specified natural source uses the simple component
\(A=\epsilon\mathbf Q[G]\simeq M_m(K)\) and its
prescribed primitive corner \(eAe=K\). These are the
coefficient and corner conventions of HD-S41 and HD-R12.
Here is the actual map that recovers the exact field;
an arbitrary extra action of \(K\) on a rational block
would not give this map with center \(K\).
For a conjugacy class \(C\) with representative \(c\),
Schur's lemma and the trace give
\[
 \begin{split}
 Z(\mathbf Q[G])&\longrightarrow Z(A)
                \longrightarrow eAe,\\
 z&\longmapsto\epsilon z\longmapsto eze,\\
 \sum_{g\in C}g&\longmapsto
        \frac{|C|\chi(c)}m\,e,
 \qquad \chi(g)=\operatorname{tr}_V(g).
 \end{split}
 \tag{F1T22.3}
\]
The first map is surjective: for a central element of the
simple direct factor, its extension by zero in the other
factors is central in \(\mathbf Q[G]\). Class sums
form a rational basis of that center. Consequently the
displayed scalar values span \(K\) over \(\mathbf Q\),
and all character values lie in \(K\) because the
original matrices do. Thus
\[
                 \mathbf Q(\chi(g):g\in G)=K.
 \tag{F1T22.4}
\]
This is precisely the exact-field premise retained here.
For the seven fields in (F1T22.1), their only roots of
unity are \(\pm1\). Absolute irreducibility therefore
makes the scalar center \(Z=Z(G)\) a subgroup of
\(\{1,-1\}\).

\paragraph{2. The actual ordinary layer and all its scalar extensions.}
By definition \(P/Z=A_n\). Its perfectness shows that
\(P'\to A_n\) is surjective, whence \(P=ZL\).
Under the present hypothesis \(L\simeq A_n\), its
center is trivial, so \(P=Z\times L\).
The actual conjugation map identifies \(G/Z\) with a
subgroup of \(\operatorname{Aut}(A_n)=S_n\) containing
\(A_n\). In particular \([G:P]\le2\). The equality
of automorphism groups for \(n\ge8\) is explicitly
used in Guralnick--Tiep, JEMS14 (2012), Proposition3.4,
printed p.626. For \(n=7\), the ATLAS record gives
\(\operatorname{Out}(A_7)=C_2\) and its actual
transposition conjugation; since the natural \(S_7\)
conjugation is faithful, it gives the same equality.

Finite averaging decomposes \(V_{\mathbf C}|_L\)
into isotypic summands; \(G\) permutes them transitively,
so there are one or two types. Each has multiplicity one.
Indeed, in the one-type case the evaluation map
\(W\otimes\operatorname{Hom}_L(W,V_{\mathbf C})
\to V_{\mathbf C}\), \(w\otimes f\mapsto f(w)\),
is an isomorphism. The group \(P\) acts only through
the first factor and scalars. A representative of
\(G/P\), when present, has the form \(A\otimes T\)
by the matrix-unit centralizer calculation. An eigenline
of \(T\) would give a proper \(G\)-stable subspace
unless the multiplicity is one. With two types, take an
irreducible copy \(W_0\) in the first isotypic space.
For \(g\notin P\), the sum \(W_0+gW_0\) is
\(G\)-stable because \(g^2\in P\); irreducibility
again forces both multiplicities to be one.

We use the following published ordinary representation
description with its explicit rational models. For each
partition \(\lambda\) of \(n\), the Young-symmetrizer
module \(S^\lambda=\mathbf Q[S_n]a_\lambda b_\lambda\)
is rational and absolutely irreducible. Its sign twist is
\(S^{\lambda'}\). If \(\lambda\ne\lambda'\),
its restriction to \(A_n\) is irreducible and fixed by
odd conjugation. If \(\lambda=\lambda'\), its
restriction is the sum of two inequivalent irreducibles
exchanged by every odd conjugation, and inducing either
one gives \(S^\lambda\). These are Fulton--Harris,
\emph{Representation Theory} (1991), Theorem4.3,
Exercise4.4 and Proposition5.1, printed pp.46--47,
52--54 and63--65; the complete Theorem4.3 and
Proposition5.1 proofs were read. The sign-twist map in
Exercise4.4 is explicit: the algebra automorphism
\(x\mapsto\eta(x)\), \(\eta(g)=\operatorname{sgn}(g)g\),
sends \(a_\lambda b_\lambda\) to
\(b_{\lambda'}a_{\lambda'}\), using transposed
tableaux. Right multiplication by \(b_{\lambda'}\)
then maps this last left ideal into
\(\mathbf Q[S_n]a_{\lambda'}b_{\lambda'}\).
It is nonzero because
\((a_{\lambda'}b_{\lambda'})^2\ne0\); composing
on the left with \(a_{\lambda'}\) verifies that
nonvanishing. Irreducibility makes it an isomorphism.
All \(S_n\) character values are integers, since
the rational model makes them rational and finite order
makes them algebraic integers.

For completeness the original extension matrices now
exclude every case except one nonreal split constituent
with \(G=P\). If \(G=P\), its single constituent
must have nonreal character by (F1T22.4), since the
central factors are only \(\pm1\). In particular it
cannot be the restriction of \(S^\lambda\) with
\(\lambda\ne\lambda'\).

Suppose instead \(G/P=C_2\). Multiply an outside
representative by an element of \(L\) so that its
actual conjugation on \(L\) is a transposition
\(\tau\). Its square centralizes \(L\), and has
trivial image in \(G/Z\), hence equals
\(z_0\in Z\). For a single constituent, odd
invariance gives the rational model
\(S^\lambda|_{A_n}\), \(\lambda\ne\lambda'\).
Choose an actual \(L\)-intertwiner to that model;
Schur's lemma gives for this outside element
\[
       \rho(g)=u\rho_\lambda(\tau),\qquad
       u^2=z_0\in\{1,-1\}.
 \tag{F1T22.5}
\]
Thus \(u\in\{1,-1,i,-i\}\). Inside \(P\)
the traces are integers up to sign. Outside they are
\(u\) times integers, up to sign, since every outside
element belongs to \(Pg\). If \(u\) is real the
entire character is rational. If \(u=\pm i\),
either all those outside traces vanish and the character
is again rational, or a nonzero one forces \(i\in K\).
Both outcomes contradict (F1T22.1) and (F1T22.4).

With two constituents, the actual \(L\)-intertwiner
identifies their sum with
\(S^\lambda|_{A_n}\), \(\lambda=\lambda'\).
Every outside matrix exchanges the two summands and
therefore has trace zero; every inside trace is the
integer \(S_n\)-character value up to sign. This also
contradicts (F1T22.4). The intertwiner here can be
constructed by the averaged Hom projections onto the
two irreducible summands, followed by their direct sum;
it does not replace the original \(G\)-action.

We have proved that the sole remaining possibility is
\[
 G=Z\times A_n,\qquad V_{\mathbf C}|_{A_n}=W,
 \qquad
 S^\lambda_{\mathbf C}|_{A_n}=W\oplus\overline W,
 \quad\lambda=\lambda',\quad \dim W=m.
 \tag{F1T22.6}
\]
Indeed the rational model is real, its two constituents
are inequivalent, and \(W\) has nonreal character.
Complex conjugation must therefore exchange them.

\paragraph{3. Exact small-support character bounds on that rational induction.}
Put \(U=S^\lambda_{\mathbf C}\), so \(\dim U=2m\),
and set
\(p_2(\lambda)=\sum_{(i,j)\in\lambda}(j-i)^2\).
The displayed \(p=3\) specialization of Lassalle,
\emph{An explicit formula for the characters of the
symmetric group}, Math. Ann.340 (2008),383--405,
Theorem4 (author version p.12), gives
\[
 r_3:=\frac{\chi_\lambda((123))}{\dim U}
   =\frac{3p_2(\lambda)-3\binom n2}
                {n(n-1)(n-2)}.
 \tag{F1T22.7}
\]
Its statement and full relevant proof chain were read.
For a self-conjugate diagram the diagonal hooks partition
the boxes. Their lengths \(h_i=2a_i+1\) are positive
odd integers with \(\sum_i h_i=n\). The contents
in the hook indexed by \(i\) are exactly
\(-a_i,\ldots,-1,0,1,\ldots,a_i\). Hence
\[
 p_2(\lambda)=2\sum_i\sum_{j=1}^{a_i}j^2
             =\frac{\sum_i h_i^3-n}{12},
 \qquad
 0\le p_2(\lambda)\le\frac{n^3-n}{12}.
 \tag{F1T22.8}
\]
Substitution into (F1T22.7) gives
\[
 -\frac3{2(n-2)}\le r_3
      \le\frac{n-5}{4(n-2)}<\frac14.
 \tag{F1T22.9}
\]

Let \(C_2\) be the sum of all transpositions in
\(\mathbf C[S_n]\). The sign self-isomorphism of
\(S^\lambda\) makes its transposition trace zero;
centrality and Schur's lemma therefore make \(C_2\)
the zero operator on \(U\). In the actual product
\(C_2^2\), fixing its first transposition leaves one
equal second transposition, \(2(n-2)\) sharing one
letter, and \(\binom{n-2}{2}\) disjoint from it.
Taking normalized traces and dividing by \(\binom n2\)
proves
\[
 0=1+2(n-2)r_3+
       \frac{(n-2)(n-3)}2r_{22},
 \qquad
 r_{22}:=\frac{\chi_\lambda((12)(34))}{\dim U}.
 \tag{F1T22.10}
\]
The lower bound in (F1T22.9) yields
\[
 r_{22}\le\frac4{(n-2)(n-3)}\le\frac15
                  \quad(n\ge7).
 \tag{F1T22.11}
\]
The exact averaging projectors for an order-three or
order-two element now give
\[
 \begin{split}
 \operatorname{rank}_U((123)-1)
   &=\frac23(2m)(1-r_3)\ge m,\\
 \operatorname{rank}_U((12)(34)-1)
   &=\frac12(2m)(1-r_{22})\ge\frac{4m}{5}.
 \end{split}
 \tag{F1T22.12}
\]
For the first equality, the two nonidentity powers of a
3-cycle are conjugate in \(S_n\), so their traces agree.

\paragraph{4. The actual induction and commutator maps bound original inertia.}
The induction used above is the literal module
\[
 \mathbf C[S_n]\otimes_{\mathbf C[A_n]}W
       \xrightarrow{\sim}U,
       \qquad a\otimes w\longmapsto\rho_\lambda(a)i(w),
 \tag{F1T22.13}
\]
where \(i:W\hookrightarrow U|_{A_n}\) is either
actual constituent embedding. The map is well defined,
nonzero and \(S_n\)-equivariant. Irreducibility of
\(U\) makes it surjective, and both sides have
dimension \(2m\), proving it is an isomorphism.
Its restriction is the two actual coset summands in
(F1T22.6).

Write an original nonidentity element as \(g=z\ell\),
with \(z\in\{1,-1\}\), \(\ell\in A_n\).
If \(\ell=1\), its moving rank is \(m\).
Otherwise choose a transposition \(\tau\) not
commuting with \(\ell\). Such a transposition exists:
commuting with every transposition would make
\(\ell\) central in \(S_n\), whose center is trivial
for \(n\ge3\). The actual commutator
\(h=[\ell,\tau]\) is the product of two distinct
transpositions. It is therefore a 3-cycle or a double
transposition.

Set \(B=z\rho_\lambda(\ell)\) and
\(T=\rho_\lambda(\tau)\) on \(U\). Because
\(z\) is real, the two summands \(W,\overline W\)
have conjugate \(B-1\) matrices, whence
\[
 \operatorname{rank}_U(B-1)
       =2\operatorname{rank}_V(g-1).
 \tag{F1T22.14}
\]
The exact matrix identity
\[
 \rho_\lambda(h)-1
   =\big((B-1)T-T(B-1)\big)B^{-1}T^{-1}
 \tag{F1T22.15}
\]
then gives, by (F1T22.12),
\[
 \frac{4m}{5}
 \le\operatorname{rank}_U(\rho_\lambda(h)-1)
 \le2\operatorname{rank}_U(B-1)
 =4\operatorname{rank}_V(g-1).
 \tag{F1T22.16}
\]
Thus every nonidentity original element has moving rank
at least \(m/5\). Applying this to all the original
inertias in (F1T22.2) proves
\[
          2m+6\ge\frac{bm}{5}\ge\frac{12m}{5},
          \qquad m\le15,
 \tag{F1T22.17}
\]
contrary to \(m\ge27\). This proves the asserted
ordinary-layer exclusion over all seven fields.

\paragraph{Original source and geometric framework.}
The center map (F1T22.3), averaged constituent maps,
actual extension matrices, Young-symmetrizer induction
and commutator identity are all maps on the original
coefficient representation or on the explicitly displayed
auxiliary induced module. The latter is used solely to
prove (F1T22.16). On the original connected cover,
coefficient local system, parabolic cohomology, selected
rational primitive projector and sixfold factor, the
comparison map is the identity. Accordingly the
parabolic cup form and restricted theta form, Poincare
biextension and its dual homomorphism, Hodge filtration,
period derivative, relative Picard and Hilbert maps,
cycle class, Grothendieck--Riemann--Roch and Adams maps
retain that same source and multiplier one. No induced
Jacobian or modified polarization is substituted in the
Euler budget (F1T22.2). The proven boundary is explicit:
the Gaussian ordinary layer and other actual central
covers are assessed by their own matrix maps.

\paragraph{Source and comparison record.}
The exact primary locators and access scopes, including
the full Young-symmetrizer and index-two proofs, the
Lassalle formula and its relevant proof chain, and the
ATLAS \(A_7\) outer automorphism record, are in
\texttt{ORDINARY\_\allowbreak ALTERNATING\_\allowbreak T22\_\allowbreak SOURCE\_\allowbreak REVIEW.md}.
The required Hexagon, arXiv and journal comparison
searches are recorded there. The new conclusion claimed
here is the exclusion with the original prescribed
nonreal corner and sixfold Euler budget; the classical
representation formulas are cited inputs.
